\chapter{Хопфовское происхождение разности сродств}
\label{ch:v10-two-reservoir-hopf-affinity}

Хопфовский цикл имеет полную силу $3\log2$ и три циклически эквивалентных
ребра. Поэтому локальная сила ребра равна
\begin{equation}
 F_e=\frac{3\log2}{3}=\log2.
\end{equation}
Для двух ванн предыдущего гейта
\begin{equation}
 a_h=-\log(1/2)=\log2,\qquad
 a_c=-\log(1/4)=2\log2,
\end{equation}
и, следовательно,
\begin{equation}
 \boxed{a_c-a_h=F_e=\log2}.
\end{equation}
Тем самым разность сродств не является новой ручной константой: она
типизированно наследуется от одного ребра хопфовского цикла.

Однако преобразование $(a_h,a_c)\mapsto(a_h+s,a_c+s)$ сохраняет разность.
Матрица условий имеет ранг/ядро $2/1$ с ядром $(1,1,0)$. Кроме того,
сопоставление $J_{\rm edge}=\kappa/3=1/66$ даёт лишь
$\kappa\Delta t=1/22$: абсолютные скорость и длительность шага
контрмасштабируются.

\begin{theorem}[Хопфовская разность сродств]
Равнорёберный цикл с силой $3\log2$ канонически поставляет разность
сродств $\log2$, необходимую двухрезервуарному току.
\end{theorem}

\begin{nogocriterion}[Разность не фиксирует две температуры]
Хопфовский цикл не выбирает общий сдвиг двух сродств и абсолютную длительность
шага. Поэтому он выводит неравновесный контраст, но не обе температуры и не
физическую скорость тока.
\end{nogocriterion}

\begin{protocol}\begin{enumerate}
\item хопфовская разность сродств: \textbf{$1/1$};
\item условное замыкание: \textbf{$8/8$};
\item общий температурный сдвиг: \textbf{$0/1$};
\item абсолютные часы: \textbf{$0/1$};
\item абсолютный масштаб: \textbf{$0/1$};
\item следующий гейт: \textbf{аудит общего температурного якоря}.
\end{enumerate}\end{protocol}

Машинный сертификат сохранён в
\path{s2t_v10_cell_birth_four_volume_induced_newton_breathing_anomaly_two_reservoir_affinity_hopf_cycle_typed_origin_gate_results.json}.